\documentclass[11pt,a4paper]{article}
\usepackage[T1]{fontenc}
\usepackage{amsmath,amssymb,amsthm}
\usepackage[margin=27mm]{geometry}
\usepackage[hidelinks]{hyperref}
\newtheorem{proposition}{Proposition}
\newcommand{\Z}{\mathbb Z}
\newcommand{\Q}{\mathbb Q}
\newcommand{\Nm}{\operatorname{N}}
\title{Magic Numbers:\\The Silver Unit in the Eighth Cyclotomic Ring}
\author{A. M. Thorn\\\small\href{https://orcid.org/0009-0008-5463-278X}{ORCID 0009-0008-5463-278X}}
\date{18 November 2026}
\begin{document}
\maketitle
\begin{abstract}
In the classical cyclotomic ring $R=\Z[m]$, with $m=e^{\pi i/4}$,
the element $M=1-m+m^2$ is a unit with an explicit integral inverse.
We give its multiplication matrix and exact forward and inverse
coordinate rules. The identity $M=m/(1+\sqrt2)$ separates an
eighth-root rotation from a silver-ratio scale. Two complex embeddings
contract and expand by reciprocal factors while preserving the quartic
norm. Starting from $1$, eight rotations return to the seed, whereas eight
$M$ steps do not. These elementary formulas provide a silver counterpart
to the golden step used in TWIST-J. They describe reversible algebra;
they do not identify $M$ with a unitary gate or derive physical dynamics.
\end{abstract}

\section{The ring and its real subring}\label{sec:ring}
Fix $m=e^{\pi i/4}$, and write
\begin{equation}\label{eq:ring}
 R=\Z[m]\simeq\Z[x]/(x^4+1),\qquad K=\Q(m),\qquad m^2=i.
\end{equation}
The polynomial $(x+1)^4+1=x^4+4x^3+6x^2+4x+2$ is
Eisenstein at $2$, so $x^4+1$ is irreducible over $\Q$.
Every element therefore has unique coordinates
\begin{equation}\label{eq:coordinates}
 \alpha=a+bm+cm^2+dm^3,\qquad (a,b,c,d)\in\Z^4.
\end{equation}
Multiplication is polynomial convolution followed by $m^4=-1$:
if the ordinary product is $\sum_{k=0}^6u_km^k$, its coordinates are
$(u_0-u_4,u_1-u_5,u_2-u_6,u_3)$. The name magic numbers refers
to this presentation of a familiar cyclotomic ring \cite{washington}.
Rank four counts integer coordinates, not physical dimensions.

Put $s=m-m^3=\sqrt2$ in the fixed embedding. Then $s^2=2$,
and complex conjugation sends $(a,b,c,d)$ to $(a,-d,-c,-b)$.
Its fixed elements have $c=0$, $d=-b$, so
\begin{equation}\label{eq:real}
 R\cap\mathbb R=\Z[s],\qquad
 \delta=1+s,\qquad \delta^{-1}=s-1,\qquad \delta^2=2\delta+1.
\end{equation}
The quadratic norm of $a+bs$ is $a^2-2b^2$; its norm in the
quartic field $K$ is $(a^2-2b^2)^2$. Thus $\delta$ has quadratic
norm $-1$ and quartic norm $1$.

The automorphisms $\sigma_r(m)=m^r$, for $r=1,3,5,7$, form the
Klein four-group: each of its three nonidentity elements has order
two, and any two distinct nonidentity elements generate the group.
There are two independent generators, not three.

\section{The integral silver step}\label{sec:step}
Define the distinguished unit
\begin{equation}\label{eq:unit}
 M=1-m+m^2=m\delta^{-1},\qquad M^{-1}=1-m^2-m^3.
\end{equation}
Indeed, $m(s-1)=m^2-m^4-m=1-m+m^2$, and
$\delta m^{-1}=(1+m-m^3)(-m^3)=1-m^2-m^3$.
These identities prove the inverse inside $R$.

In the ordered basis $(1,m,m^2,m^3)$, use column vectors and let
$A$ denote multiplication by $m$. Then
\begin{equation}\label{eq:matrix}
 A=\begin{pmatrix}0&0&0&-1\\1&0&0&0\\0&1&0&0\\0&0&1&0\end{pmatrix},
 \qquad B=I-A+A^2=
 \begin{pmatrix}1&0&-1&1\\-1&1&0&-1\\1&-1&1&0\\0&1&-1&1\end{pmatrix}.
\end{equation}
The signed shift $A$ satisfies $A^4=-I$ and $A^8=I$.
The matrix $B$ represents multiplication by $M$.

\begin{proposition}\label{prop:step}
Multiplication by $M$ and by $M^{-1}$ gives the mutually inverse rules
\begin{align}
 (a,b,c,d)&\longmapsto(a-c+d,\ b-a-d,\ c-b+a,\ d-c+b),\label{eq:forward}\\
 (A,B,C,D)&\longmapsto(A+B+C,\ B+C+D,\ -A+C+D,\ -A-B+D).\label{eq:inverse}
\end{align}
Both preserve integer coordinates and the quartic field norm.
Moreover,
\begin{equation}\label{eq:charpoly}
 \det B=1,\qquad \det(xI-B)=x^4-4x^3+8x^2-4x+1.
\end{equation}
\end{proposition}
\begin{proof}
Reduce the four products of $M$ with the basis elements to obtain
\eqref{eq:matrix} and \eqref{eq:forward}. Multiplying instead by
$1-m^2-m^3$ gives \eqref{eq:inverse}; substitution verifies both
compositions. The characteristic polynomial and norm preservation
follow from the embeddings below. The coordinate operations use
only additions and subtractions. This does not imply bounded
coefficient sizes or constant bit cost under iteration.
\end{proof}

\section{Two embeddings and one norm}\label{sec:embeddings}
Choose $\sigma_1$ and $\sigma_3$, one from each complex-conjugate
pair. Under $\sigma_3$ the element $s$ becomes $-\sqrt2$ and
$\delta$ becomes $1-\sqrt2=-\delta^{-1}$, where the positive real
number $\delta=1+\sqrt2$ denotes its value in the fixed embedding.
Consequently,
\begin{equation}\label{eq:embeddings}
 \sigma_1(M)=\delta^{-1}e^{\pi i/4},\qquad
 \sigma_3(M)=\delta e^{-\pi i/4}.
\end{equation}
For the second equality, $m^3/(1-\sqrt2)=-\delta m^3$
has argument $-\pi/4$. Thus the same exact step contracts one
complex plane by $\delta^{-1}$ and expands the other by $\delta$.

For any $\alpha\in K$, define
\begin{equation}\label{eq:norm}
 \Nm_{K/\Q}(\alpha)=\prod_{r=1,3,5,7}\sigma_r(\alpha)
 =|\sigma_1(\alpha)|^2|\sigma_3(\alpha)|^2.
\end{equation}
The evaluation matrix $(m^{rk})$, with $r=1,3,5,7$ and
$k=0,1,2,3$, is Vandermonde at four distinct roots. It is invertible
and diagonalizes multiplication by $\alpha$ with eigenvalues
$\sigma_r(\alpha)$. Hence the field norm is the determinant of
the multiplication matrix, an integer when $\alpha\in R$.
The real-linear extension of $(\sigma_1,\sigma_3)$ identifies
$K\otimes_{\Q}\mathbb R$ with $\mathbb C^2$.
The determinant of multiplication by $M$ is
$\delta^{-2}\delta^2=1$, and multiplicativity proves norm preservation.

The two conjugate pairs in \eqref{eq:embeddings} give
\[
 \det(xI-B)=
 (x^2-(2-\sqrt2)x+3-2\sqrt2)
 (x^2-(2+\sqrt2)x+3+2\sqrt2).
\]
Expansion proves \eqref{eq:charpoly}. No eigenvalue has modulus
one, so $B$ is hyperbolic. In particular $M$ has infinite order:
$|\sigma_1(M^n)|=\delta^{-n}$ equals $1$ only when $n=0$.
An integral unit therefore need not preserve length in either complex
plane or in the ordinary coefficient metric.

\section{Eight steps and the boundary of the result}\label{sec:scope}
The rotation and the silver step have the same standard phase per
step, but different scales. From \eqref{eq:unit},
\begin{equation}\label{eq:eight}
 m^8=1,\qquad M^8=\delta^{-8}=577-408s=577-408m+408m^3.
\end{equation}
For example, $\delta^{-2}=3-2s$, $\delta^{-4}=17-12s$, and
squaring again gives \eqref{eq:eight}. Starting at $1$, eight pure
rotations return to $(1,0,0,0)$; eight silver steps give
$(577,-408,0,408)$. Phase closure is not orbit closure.

The superficially similar element $1+m^2=1+i=sm$ is not a unit.
Its norm is $4$: in the two selected embeddings its squared moduli
are both $2$. This contrasts with $M$, whose norm is $1$.
The contrast concerns invertibility over the integral ring, not
physical irreversibility or measurement.

The algebra gives a silver counterpart of the golden relation
$\varphi J=j$ used in TWIST-J, where $j=e^{2\pi i/5}$,
$J=1+j^2$ and $\varphi=(1+\sqrt5)/2$: here $\delta M=m$.
It does not assign forces, select a native update law or identify
integer rank with spacetime dimension. Nor is $M$ a unitary quantum
gate: its complex modulus in the fixed embedding is $\delta^{-1}$,
not $1$. The phase matrix $\operatorname{diag}(1,m)$ is unitary;
it is a different mathematical object. A coefficient ring and a
level of the Clifford hierarchy are likewise different kinds of object.
No quantum synthesis or universality theorem is claimed here.

\paragraph{Verification.}
The accompanying \texttt{verify\_algebra.py} checks the identities
by integer polynomial reduction, basis tests and determinant expansion.
These finite calculations check the displayed algebra; the proofs
establish the stated general results.

\paragraph{Edition.}
An earlier edition appeared as \emph{Magic Numbers and
Four-Dimensional Integer M Algebra}
(\href{https://doi.org/10.5281/zenodo.20560040}{Zenodo 20560040}).

\begin{thebibliography}{9}
\bibitem{washington} L. C. Washington, \emph{Introduction to Cyclotomic
Fields}, second edition, Springer, 1997.
\href{https://doi.org/10.1007/978-1-4612-1934-7}{Publisher DOI}.
\end{thebibliography}
\end{document}
